% Shared academic preamble. White page, black text, serif body.
\documentclass[a4paper,11pt,fontset=none]{ctexart}
\setmainfont{Noto Serif}
\setCJKmainfont{Noto Serif CJK SC}[BoldFont={Noto Serif CJK SC Bold}]
\setCJKsansfont{Noto Sans CJK SC}[BoldFont={Noto Sans CJK SC Bold}]
\setCJKmonofont{Noto Sans CJK SC}
\usepackage{geometry}
\geometry{a4paper,margin=2.3cm,headheight=14pt}
\usepackage{booktabs}
\usepackage{longtable}
\usepackage{array}
\usepackage{tabularx}
\usepackage{tikz}
\usetikzlibrary{arrows.meta,positioning}
\usepackage{float}
\usepackage{caption}
\captionsetup{font=small,labelfont=bf,skip=6pt}
\usepackage{fancyhdr}
\usepackage{setspace}
\setstretch{1.15}
\usepackage{enumitem}
\setlist[itemize]{itemsep=0.2em,topsep=0.35em,leftmargin=1.6em}
\setlist[enumerate]{itemsep=0.2em,topsep=0.35em,leftmargin=1.6em}
\usepackage{amsmath}
\usepackage{xurl}
\usepackage[hidelinks]{hyperref}
\pagestyle{fancy}
\fancyhf{}
\fancyhead[L]{\small 盼盼法式小面包}
\fancyhead[R]{}
\fancyfoot[C]{\thepage}
\renewcommand{\headrulewidth}{0.4pt}
\setcounter{secnumdepth}{3}
\setcounter{tocdepth}{2}
\emergencystretch=2em
\Urlmuskip=0mu plus 1mu
\tikzset{
  layer/.style={draw, align=flush center, font=\small, inner sep=5pt, text width=3.35cm, minimum height=1.05cm},
  wide/.style={draw, align=flush center, font=\small, inner sep=6pt, text width=14.2cm},
  flow/.style={-{Stealth[length=2.2mm]}, semithick},
  abox/.style={draw, fill=white, align=flush center, font=\scriptsize, inner sep=2.5pt,
    text width=27mm, minimum height=11mm},
  sbox/.style={draw, dashed, fill=white, align=flush center, font=\scriptsize, inner sep=2.5pt,
    text width=24mm, minimum height=11mm},
  alab/.style={font=\scriptsize\bfseries, align=center, text width=12mm},
  aflow/.style={-{Stealth[length=1.6mm]}, semithick},
  aside/.style={-{Stealth[length=1.6mm]}, densely dashed, semithick}
}
\newcommand{\contesttitlepage}[6]{%
  \begin{titlepage}
  \thispagestyle{empty}
  \centering
  \vspace*{1.5cm}
  {\large 首届江苏省“人工智能+”创新应用大赛\par}
  \vspace{0.55cm}
  {\large #1\par}
  \vspace{1.35cm}
  {\rule{0.42\textwidth}{0.5pt}\par}
  \vspace{0.75cm}
  {\Huge\bfseries #2\par}
  \vspace{0.4cm}
  {\Large #3\par}
  \vspace{0.95cm}
  {\LARGE #4\par}
  \vspace{0.5cm}
  {\large #5\par}
  \vspace{0.75cm}
  {\rule{0.42\textwidth}{0.5pt}\par}
  \vfill
  {\large #6\par}
  \vspace*{1.1cm}
  \end{titlepage}
}

\hypersetup{pdftitle={FluVax-Evidence 项目计划书},pdfauthor={盼盼法式小面包}}
\begin{document}
\contesttitlepage
  {“人工智能+”医疗健康赛道\\赛题 AI-HEALTH-Q40\\流感病毒免疫原辅助设计智能体研发和应用}
  {FluVax-Evidence}
  {盼盼法式小面包}
  {项目计划书}
  {已发表公开记录的离线证据审阅}
  {创新组\quad 报名编号 REG-2026-000525\\命题单位：宿迁市楚润数据集团有限公司\\2026 年 11 月}

\begin{abstract}
FluVax-Evidence 把已发表公开记录收成可签批的问答和方案档案对照。7 道封闭题的 Recall@5 为 1.0，NDCG@5 为 0.9285714285714286；四份档案与金标一致；已通过原文与版本锁的语句为 59/59。人工检索 SOP 未测。团队盼盼法式小面包，创新组。
\end{abstract}

\section{问题与范围}
流感免疫原相关资料分散在文献、核酸库、结构库和表位库中，版本和出处需要留在每一条结论上。赛题要求软件原型完成检索、关联、既有方案比较和报告，并由专业人员复核。\footnote{范围：软件原型与公开记录离线评测。湿实验、序列设计和疗效结论不在本阶段交付内。人工检索 SOP 未测；草案中的耗时与召回目标未测。}

命题单位是宿迁市楚润数据集团有限公司。团队报名编号 REG-2026-000525。

\section{方案}
图~\ref{fig:arch} 为四层结构。数据层锁定 24 条公开记录（版本 fluevidence-v1.1）：文献题名、核酸库定义行、结构库题名、表位摘要和蛋白质条目首行。算法层合并词频检索、384 维稠密向量和 1 至 3 跳图检索，再用 362 参数的重排器和 387 参数的忠实度判别器。服务层核对逐字引文与版本，并比较四份已发表方案档案。交互层给出问答、对照表和签批；签批前报告保持草稿。默认把选中记录的原文装成引文。外部大模型本次未接入。

\begin{figure}[H]
\centering
\begin{tikzpicture}[x=1mm,y=1mm]
  \draw[black!40, fill=black!3] (18,71) rectangle (112,85);
  \draw[black!40, fill=black!3] (18,49) rectangle (112,63);
  \draw[black!40, fill=black!3] (18,27) rectangle (112,41);
  \draw[black!40, fill=black!3] (18,5) rectangle (112,19);
  \node[alab] at (8,78) {交互};
  \node[alab] at (8,56) {服务};
  \node[alab] at (8,34) {算法};
  \node[alab] at (8,12) {数据};
  \node[abox] at (35,78) {循证问答};
  \node[abox] at (67,78) {四份档案对照};
  \node[abox] at (99,78) {专家签批};
  \node[abox] at (35,56) {发表与入库\\时间汇总};
  \node[abox] at (67,56) {逐字与版本核验};
  \node[abox] at (99,56) {来源与等级比较};
  \node[abox] at (35,34) {词频与精确标识};
  \node[abox] at (67,34) {稠密向量与图};
  \node[abox] at (99,34) {重排与忠实度};
  \node[abox] at (35,12) {24 条公开记录};
  \node[abox] at (67,12) {版本锁};
  \node[abox] at (99,12) {四份文献档案};
  \node[sbox] at (136,34) {离线原文组装};
  \node[sbox] at (136,56) {疗效语句检查};
  \draw[aflow] (67,19) -- (67,27);
  \draw[aflow] (67,41) -- (67,49);
  \draw[aflow] (67,63) -- (67,71);
  \draw[aside] (124,34) -- (112,34);
  \draw[aside] (124,56) -- (112,56);
  \node[font=\scriptsize, anchor=west] at (18,0) {实线：主数据流。虚线框：旁路。};
\end{tikzpicture}
\caption{FluVax-Evidence 分层架构。}
\label{fig:arch}
\end{figure}

证据等级按公开记录的类型标注：结构与表位实验登录 6 条，核酸库元数据 2 条，文献与审校蛋白条目 16 条，计算预测 0 条。标识核验日为 2026-11-28。\footnote{待核实：本次未重新请求核酸库、结构库、表位库和蛋白质库接口。}

重排器在合成测试集上“相关”类 F1 为 0，宏平均 F1 为 0.42857142857142855，NDCG@3 为 0.9692441461309548，因此排序用相关性概率。忠实度判别器在合成留出集上三类宏平均 F1 为 1.0。两套权重都在通用合成句上训练。

\section{实测}
任务一 7 题、取前 5 条：Recall@5 $=1.0$，NDCG@5 $=0.9285714285714286$，中位 18.16363 毫秒。第 6 题的目标记录位于第三位，该题 NDCG@5 为 0.5。任务二比较 4 份档案，16.912775 毫秒完成比较，11.485772 毫秒完成签批，采纳 23 条、标疑 1 条，4 份均与金标一致。已采纳语句与版本锁一致者为 59/59。不读原文的对照草稿 3 条全部被拒。疗效承诺探针 1 条被拒。方法学夹具另有 4 题，Recall@3 $=1.0$、NDCG@3 $=0.875$，与上述 7 题分开。逐题与逐档案表见佐证。

\section{创新与应用}
每条结论带回来源、数据版本和证据等级。档案对照同时给出来源覆盖、版本是否一致、公开标识交叉和等级分布。专家签批前，报告保持草稿。疗效承诺语句在门禁处被拒绝。

使用方式是院所或数据机构的内网资料整理：问答和四份方案的对照表供专业人员签批。仓库没有签约、报价和到户部署，金额不估算，商务接触标为待核实。运行记录中的环境为 Windows 与 Python 3.13，评测可对本地快照断网进行。

团队为盼盼法式小面包。初赛材料截止 2026-11-30。

\begin{thebibliography}{9}
\bibitem{challenge} 赛题 AI-HEALTH-Q40。命题单位：宿迁市楚润数据集团有限公司。
\bibitem{bench} FluEvidenceBench-v1.1，24 条公开记录。评测数字见佐证。
\end{thebibliography}
\end{document}
